\documentclass[10pt,a4paper]{amsart}
\usepackage{amsmath,amssymb,amsthm,graphicx}
\usepackage{verbatim}
\newenvironment{reference}{\comment}{\endcomment}
%\newenvironment{reference}{}{}
\newenvironment{slogan}{\comment}{\endcomment}
\newenvironment{history}{\comment}{\endcomment}

% For commutative diagrams we use Xy-pic
\usepackage[all]{xy}

% We use 2cell for 2-commutative diagrams.
\xyoption{2cell}
% Theorem environments.
%
\theoremstyle{plain}
\newtheorem{theorem}[subsection]{Theorem}
\newtheorem{proposition}[subsection]{Proposition}
\newtheorem{lemma}[subsection]{Lemma}

\theoremstyle{definition}
\newtheorem{definition}[subsection]{Definition}
\newtheorem{example}[subsection]{Example}
\newtheorem{exercise}[subsection]{Exercise}
\newtheorem{situation}[subsection]{Situation}

\theoremstyle{remark}
\newtheorem{remark}[subsection]{Remark}
\newtheorem{remarks}[subsection]{Remarks}

\numberwithin{equation}{subsection}

% Macros
%
\def\lim{\mathop{\mathrm{lim}}\nolimits}
\def\colim{\mathop{\mathrm{colim}}\nolimits}
\def\Spec{\mathop{\mathrm{Spec}}}
\def\Hom{\mathop{\mathrm{Hom}}\nolimits}
\def\Ext{\mathop{\mathrm{Ext}}\nolimits}
\def\SheafHom{\mathop{\mathcal{H}\!\mathit{om}}\nolimits}
\def\SheafExt{\mathop{\mathcal{E}\!\mathit{xt}}\nolimits}
\def\Sch{\mathit{Sch}}
\def\Mor{\mathop{\mathrm{Mor}}\nolimits}
\def\Ob{\mathop{\mathrm{Ob}}\nolimits}
\def\Sh{\mathop{\mathit{Sh}}\nolimits}
\def\NL{\mathop{N\!L}\nolimits}
\def\CH{\mathop{\mathrm{CH}}\nolimits}
\def\proetale{{pro\text{-}\acute{e}tale}}
\def\etale{{\acute{e}tale}}
\def\QCoh{\mathit{QCoh}}
\def\Ker{\mathop{\mathrm{Ker}}}
\def\Im{\mathop{\mathrm{Im}}}
\def\Coker{\mathop{\mathrm{Coker}}}
\def\Coim{\mathop{\mathrm{Coim}}}

% Boxtimes
%
\DeclareMathSymbol{\boxtimes}{\mathbin}{AMSa}{"02}

%
% Macros for moduli stacks/spaces
%
\def\QCohstack{\mathcal{QC}\!\mathit{oh}}
\def\Cohstack{\mathcal{C}\!\mathit{oh}}
\def\Spacesstack{\mathcal{S}\!\mathit{paces}}
\def\Quotfunctor{\mathrm{Quot}}
\def\Hilbfunctor{\mathrm{Hilb}}
\def\Curvesstack{\mathcal{C}\!\mathit{urves}}
\def\Polarizedstack{\mathcal{P}\!\mathit{olarized}}
\def\Complexesstack{\mathcal{C}\!\mathit{omplexes}}
% \Pic is the operator that assigns to X its picard group, usage \Pic(X)
% \Picardstack_{X/B} denotes the Picard stack of X over B
% \Picardfunctor_{X/B} denotes the Picard functor of X over B
\def\Pic{\mathop{\mathrm{Pic}}\nolimits}
\def\Picardstack{\mathcal{P}\!\mathit{ic}}
\def\Picardfunctor{\mathrm{Pic}}
\def\Deformationcategory{\mathcal{D}\!\mathit{ef}}
\usepackage{cleveref}
\crefname{originalreference}{original reference}{original references}
\Crefname{originalreference}{Original reference}{Original references}
\makeatletter
\newlabel{lemma-finite-fibres-presentation}{{03IJ}{source}{}{sourcepointer-1}{}}
\newlabel{lemma-finite-fibres-presentation@cref}{{[originalreference][0][]03IJ}{[1][0][]source}{}{}{}}
\makeatother
\crefname{originalreference}{original reference}{original references}
\Crefname{originalreference}{Original reference}{Original references}
\makeatletter
\newlabel{lemma-points-presentation}{{03BW}{source}{}{sourcepointer-2}{}}
\newlabel{lemma-points-presentation@cref}{{[originalreference][0][]03BW}{[1][0][]source}{}{}{}}
\makeatother
\crefname{originalreference}{original reference}{original references}
\Crefname{originalreference}{Original reference}{Original references}
\makeatletter
\newlabel{lemma-quasi-separated-quasi-compact-pieces}{{03W7}{source}{}{sourcepointer-3}{}}
\newlabel{lemma-quasi-separated-quasi-compact-pieces@cref}{{[originalreference][0][]03W7}{[1][0][]source}{}{}{}}
\makeatother
\crefname{originalreference}{original reference}{original references}
\Crefname{originalreference}{Original reference}{Original references}
\makeatletter
\newlabel{lemma-reduced-closed-subspace}{{03IQ}{source}{}{sourcepointer-4}{}}
\newlabel{lemma-reduced-closed-subspace@cref}{{[originalreference][0][]03IQ}{[1][0][]source}{}{}{}}
\makeatother
\crefname{originalreference}{original reference}{original references}
\Crefname{originalreference}{Original reference}{Original references}
\makeatletter
\newlabel{lemma-topology-points}{{03BX}{source}{}{sourcepointer-5}{}}
\newlabel{lemma-topology-points@cref}{{[originalreference][0][]03BX}{[1][0][]source}{}{}{}}
\makeatother
\crefname{originallemma}{original lemma}{original lemmas}
\Crefname{originallemma}{Original lemma}{Original lemmas}
\makeatletter
\newlabel{morphisms-lemma-etale-flat}{{02GS}{source}{}{sourcepointer-6}{}}
\newlabel{morphisms-lemma-etale-flat@cref}{{[originallemma][0][]02GS}{[1][0][]source}{}{}{}}
\makeatother
\crefname{originallemma}{original lemma}{original lemmas}
\Crefname{originallemma}{Original lemma}{Original lemmas}
\makeatletter
\newlabel{morphisms-lemma-generalizations-lift-flat}{{03HV}{source}{}{sourcepointer-7}{}}
\newlabel{morphisms-lemma-generalizations-lift-flat@cref}{{[originallemma][0][]03HV}{[1][0][]source}{}{}{}}
\makeatother
\crefname{originalreference}{original reference}{original references}
\Crefname{originalreference}{Original reference}{Original references}
\makeatletter
\newlabel{proposition-locally-quasi-separated-open-dense-scheme}{{06NH}{source}{}{sourcepointer-8}{}}
\newlabel{proposition-locally-quasi-separated-open-dense-scheme@cref}{{[originalreference][0][]06NH}{[1][0][]source}{}{}{}}
\makeatother
\crefname{originallemma}{original lemma}{original lemmas}
\Crefname{originallemma}{Original lemma}{Original lemmas}
\makeatletter
\newlabel{schemes-lemma-scheme-sober}{{01IS}{source}{}{sourcepointer-9}{}}
\newlabel{schemes-lemma-scheme-sober@cref}{{[originallemma][0][]01IS}{[1][0][]source}{}{}{}}
\makeatother
\crefname{originaldefinition}{original definition}{original definitions}
\Crefname{originaldefinition}{Original definition}{Original definitions}
\makeatletter
\newlabel{topology-definition-generic-point}{{004X}{source}{}{sourcepointer-10}{}}
\newlabel{topology-definition-generic-point@cref}{{[originaldefinition][0][]004X}{[1][0][]source}{}{}{}}
\makeatother
\crefname{originallemma}{original lemma}{original lemmas}
\Crefname{originallemma}{Original lemma}{Original lemmas}
\makeatletter
\newlabel{topology-lemma-quotient-kolmogorov}{{06NA}{source}{}{sourcepointer-11}{}}
\newlabel{topology-lemma-quotient-kolmogorov@cref}{{[originallemma][0][]06NA}{[1][0][]source}{}{}{}}
\makeatother
\crefname{originallemma}{original lemma}{original lemmas}
\Crefname{originallemma}{Original lemma}{Original lemmas}
\makeatletter
\newlabel{topology-lemma-sober-local}{{06N9}{source}{}{sourcepointer-12}{}}
\newlabel{topology-lemma-sober-local@cref}{{[originallemma][0][]06N9}{[1][0][]source}{}{}{}}
\makeatother
\title{The underlying topological space of a Zariski locally quasi-separated algebraic space is sober}
\author{The Stacks Project Authors}
\date{stacks-project:spaces-properties.tex:local-snapshot-20261012; GNU FDL 1.2 or later}
\begin{document}
\maketitle
\noindent\textit{Editorial note.} All delivered mathematical text is the author's own, copied from the pinned Stacks Project chapter file spaces-properties.tex; no statement or proof step was rewritten, reordered or removed.\par\medskip
\noindent\textit{Editorial note.} The counted proof is the author's complete argument: the reduction to a single affine etale chart with the finite fibre Kolmogorov criterion, the production of a closed subspace with an open dense subscheme for an irreducible closed subset, and the closing use of soberness of schemes to obtain a generic point.\par\medskip
\noindent\textit{Editorial note.} The label lemma-quasi-separated-sober names the counted statement, and the proof environment immediately following it in the source is delivered verbatim as its proof.\par\medskip
\noindent\textit{Editorial note.} The counted statement is carried by the author's own lemma environment in the chapter source.\par\medskip
\noindent\textit{Editorial note.} Pre-open gate: python operations/agent-authoring/tools/precard\_gate.py --source spaces-properties.tex --label lemma-quasi-separated-sober --start 2079 --end 2125 --local-id 28046 returned RESULT=OK (L28046 free, occupancy FREE, public interval clear). Candidate taken from sources/stacks-spaces-properties/candidates-halfA.json.\par\medskip
\noindent\textit{Editorial note.} The statement and proof are copied from the pinned chapter file; the cited inputs are kept as external pointers with their canonical Stacks tags. The proof carries a footnote listing the auxiliary facts used to verify the hypotheses of the Kolmogorov criterion; that footnote is delivered verbatim.\par\medskip
\noindent\textit{Allowed standard prerequisites (editorial).}\begin{itemize}
\item Local nature of soberness. Topology, the lemma that soberness of a topological space can be checked on the members of an open covering. Supplies the reduction to a single piece of a Zariski cover, using the characterization of Zariski local quasi-separatedness by covers with affine etale maps.
\item Kolmogorov quotient criterion for a finite fibre equivalence relation. Topology, the lemma giving conditions under which the quotient of a space by an equivalence relation with finite fibres is Kolmogorov. Supplies the conclusion that the quotient space is Kolmogorov from the finiteness of the fibres of the two projections.
\item Existence of an open dense subscheme of a locally quasi-separated space. Same chapter, the proposition producing an open dense subspace which is a scheme. Supplies the open dense subscheme Z' of the closed subspace Z.
\item Soberness of schemes and finiteness of the fibres of the projections. Schemes, the lemma that schemes are sober; Same chapter, the lemma that the fibres of the projections of the presentation are finite. Supply the generic point of the irreducible scheme Z' and the finiteness hypothesis of the Kolmogorov criterion.
\end{itemize}\medskip
\ifdefined\hypertarget\hypertarget{sourcepointer-1}{}\fi
\noindent\textit{Source reference.} Original reference 03IJ: sources/stacks-spaces-properties/pinned-original-source.tex, source lines 1155--1155. Prerequisite pointer invoked at proof line(s) 2094; it names a standard fact of the same chapter that the counted argument uses, and every step of the counted proof built on it is retained verbatim in the delivered body.\par\medskip
\ifdefined\hypertarget\hypertarget{sourcepointer-2}{}\fi
\noindent\textit{Source reference.} Original reference 03BW: sources/stacks-spaces-properties/pinned-original-source.tex, source lines 569--569. Prerequisite pointer invoked at proof line(s) 2103; it names a standard fact of the same chapter that the counted argument uses, and every step of the counted proof built on it is retained verbatim in the delivered body.\par\medskip
\ifdefined\hypertarget\hypertarget{sourcepointer-3}{}\fi
\noindent\textit{Source reference.} Original reference 03W7: sources/stacks-spaces-properties/pinned-original-source.tex, source lines 1065--1065. Prerequisite pointer invoked at proof line(s) 2090; it names a standard fact of the same chapter that the counted argument uses, and every step of the counted proof built on it is retained verbatim in the delivered body.\par\medskip
\ifdefined\hypertarget\hypertarget{sourcepointer-4}{}\fi
\noindent\textit{Source reference.} Original reference 03IQ: sources/stacks-spaces-properties/pinned-original-source.tex, source lines 1684--1684. Prerequisite pointer invoked at proof line(s) 2110; it names a standard fact of the same chapter that the counted argument uses, and every step of the counted proof built on it is retained verbatim in the delivered body.\par\medskip
\ifdefined\hypertarget\hypertarget{sourcepointer-5}{}\fi
\noindent\textit{Source reference.} Original reference 03BX: sources/stacks-spaces-properties/pinned-original-source.tex, source lines 596--596. Prerequisite pointer invoked at proof line(s) 2105; it names a standard fact of the same chapter that the counted argument uses, and every step of the counted proof built on it is retained verbatim in the delivered body.\par\medskip
\ifdefined\hypertarget\hypertarget{sourcepointer-6}{}\fi
\noindent\textit{Source reference.} Original lemma 02GS: sources/stacks-chow/tags.txt, source lines 3239--3239. Prerequisite pointer invoked at proof line(s) 2100; it names a standard fact of a sibling chapter that the counted argument uses, and every step of the counted proof built on it is retained verbatim in the delivered body.\par\medskip
\ifdefined\hypertarget\hypertarget{sourcepointer-7}{}\fi
\noindent\textit{Source reference.} Original lemma 03HV: sources/stacks-chow/tags.txt, source lines 4581--4581. Prerequisite pointer invoked at proof line(s) 2101; it names a standard fact of a sibling chapter that the counted argument uses, and every step of the counted proof built on it is retained verbatim in the delivered body.\par\medskip
\ifdefined\hypertarget\hypertarget{sourcepointer-8}{}\fi
\noindent\textit{Source reference.} Original reference 06NH: sources/stacks-spaces-properties/pinned-original-source.tex, source lines 1864--1864. Prerequisite pointer invoked at proof line(s) 2117; it names a standard fact of the same chapter that the counted argument uses, and every step of the counted proof built on it is retained verbatim in the delivered body.\par\medskip
\ifdefined\hypertarget\hypertarget{sourcepointer-9}{}\fi
\noindent\textit{Source reference.} Original lemma 01IS: sources/stacks-chow/tags.txt, source lines 1989--1989. Prerequisite pointer invoked at proof line(s) 2099; it names a standard fact of a sibling chapter that the counted argument uses, and every step of the counted proof built on it is retained verbatim in the delivered body.\par\medskip
\ifdefined\hypertarget\hypertarget{sourcepointer-10}{}\fi
\noindent\textit{Source reference.} Original definition 004X: sources/stacks-chow/tags.txt, source lines 189--189. Prerequisite pointer invoked at proof line(s) 2083; it names a standard fact of a sibling chapter that the counted argument uses, and every step of the counted proof built on it is retained verbatim in the delivered body.\par\medskip
\ifdefined\hypertarget\hypertarget{sourcepointer-11}{}\fi
\noindent\textit{Source reference.} Original lemma 06NA: sources/stacks-chow/tags.txt, source lines 8529--8529. Prerequisite pointer invoked at proof line(s) 2096; it names a standard fact of a sibling chapter that the counted argument uses, and every step of the counted proof built on it is retained verbatim in the delivered body.\par\medskip
\ifdefined\hypertarget\hypertarget{sourcepointer-12}{}\fi
\noindent\textit{Source reference.} Original lemma 06N9: sources/stacks-chow/tags.txt, source lines 8528--8528. Prerequisite pointer invoked at proof line(s) 2088; it names a standard fact of a sibling chapter that the counted argument uses, and every step of the counted proof built on it is retained verbatim in the delivered body.\par\medskip
\begin{lemma}
\label{lemma-quasi-separated-sober}
Let $S$ be a scheme. Let $X$ be a Zariski locally quasi-separated
algebraic space over $S$. Then the topological space $|X|$ is sober (see
Topology, Definition \ref{topology-definition-generic-point}).
\end{lemma}

\begin{proof}
Combining
Topology, Lemma \ref{topology-lemma-sober-local}
and
Lemma \ref{lemma-quasi-separated-quasi-compact-pieces}
we see that we may assume that there exists an affine scheme $U$
and a surjective, quasi-compact, \'etale morphism $U \to X$.
Set $R = U \times_X U$ with projection maps $s, t : R \to U$. Applying
Lemma \ref{lemma-finite-fibres-presentation}
we see that the fibres of $s, t$ are finite. It follows all the assumptions of
Topology, Lemma \ref{topology-lemma-quotient-kolmogorov}
are met, and we conclude that $|X|$ is Kolmogorov\footnote{
Actually we use here also
Schemes, Lemma \ref{schemes-lemma-scheme-sober} (soberness schemes),
Morphisms, Lemmas \ref{morphisms-lemma-etale-flat}
and \ref{morphisms-lemma-generalizations-lift-flat} (generalizations
lift along \'etale morphisms),
Lemma \ref{lemma-points-presentation} (points on an algebraic space in
terms of a presentation), and
Lemma \ref{lemma-topology-points} (openness quotient map).}.

\medskip\noindent
It remains to show that every irreducible closed subset
$T \subset |X|$ has a generic point. By
Lemma \ref{lemma-reduced-closed-subspace}
there exists a closed subspace $Z \subset X$ with $|Z| = |T|$.
Note that $U \times_X Z \to Z$ is a quasi-compact, surjective, \'etale
morphism from an affine scheme to $Z$, hence $Z$ is Zariski locally
quasi-separated by
Lemma \ref{lemma-quasi-separated-quasi-compact-pieces}.
By
Proposition \ref{proposition-locally-quasi-separated-open-dense-scheme}
we see that there exists an open dense subspace $Z' \subset Z$
which is a scheme. This means that $|Z'| \subset T$ is open dense.
Hence the topological space $|Z'|$ is irreducible, which means that
$Z'$ is an irreducible scheme. By
Schemes, Lemma \ref{schemes-lemma-scheme-sober}
we conclude that $|Z'|$ is the closure of a single point
$\eta \in |Z'| \subset T$ and hence also $T = \overline{\{\eta\}}$, and we win.
\end{proof}

\end{document}
